\documentclass{article}

% NeurIPS 2026 template.
% "preprint" keeps author names visible and removes review line numbers.
% Remove "preprint" if your instructor wants the anonymous submission style.
\usepackage[preprint]{neurips_2026}

\usepackage[utf8]{inputenc}
\usepackage[T1]{fontenc}
\usepackage{amsmath}
\usepackage{amssymb}
\usepackage{amsfonts}
\usepackage{graphicx}
\usepackage{booktabs}
\usepackage{microtype}
\usepackage{url}
\usepackage{hyperref}

\title{Kernel Density Estimation via Diffusion}

\author{
  Erik Quddus
  \And
  Badr Eddin Amaya
}

\begin{document}

\maketitle


% ============================================================
% (a) MAIN ALGORITHM, FORMULATION, AND THESIS
% ============================================================

\section{Method and main idea}

Given observations
\[
X_1,\ldots,X_N \overset{\mathrm{iid}}{\sim} f,
\]
the goal is to estimate the unknown density $f$. The ordinary Gaussian
kernel density estimator (KDE) is
\[
\widehat f(x;t)
=
\frac{1}{N}
\sum_{i=1}^{N}
\phi(x,X_i;t),
\qquad
\phi(x,y;t)
=
\frac{1}{\sqrt{2\pi t}}
\exp\left(
-\frac{(x-y)^2}{2t}
\right),
\]
where the conventional bandwidth is $h=\sqrt{t}$.

A key observation of Botev, Grotowski, and Kroese
\citep{botev2010} is that the Gaussian KDE is also the solution of the
heat equation
\[
\frac{\partial \widehat f}{\partial t}
=
\frac{1}{2}
\frac{\partial^2 \widehat f}{\partial x^2},
\qquad
\widehat f(x;0)
=
\frac{1}{N}
\sum_{i=1}^{N}\delta(x-X_i).
\]
Thus, $t$ can be interpreted as diffusion time: increasing $t$
corresponds to increasing the amount of smoothing.

The paper replaces the ordinary heat equation by the more general
diffusion equation
\[
\frac{\partial g(x;t)}{\partial t}
=
Lg(x;t),
\]
with
\[
Lh(x)
=
\frac{1}{2}
\frac{d}{dx}
\left[
a(x)
\frac{d}{dx}
\left(
\frac{h(x)}{p(x)}
\right)
\right].
\]
Here $p(x)$ is a pilot density and $a(x)$ determines the local diffusion
behaviour. The family
\[
a(x)=p(x)^\alpha,
\qquad 0\leq\alpha\leq1,
\]
includes several forms of adaptive smoothing. Following the numerical
experiments in the paper, we use $\alpha=1$, so that
\[
a(x)=p(x).
\]

The main thesis of the paper is that interpreting KDE as diffusion
provides a flexible adaptive estimator that can use information from a
pilot density, naturally handle boundaries, and combine this with a
fully data-driven plug-in bandwidth selector that avoids the normal
reference rules used by earlier plug-in methods.


\subsection{Algorithm 1: pilot bandwidth selection}

The first stage is the Improved Sheather--Jones (ISJ) procedure. Its
purpose is to obtain a bandwidth for a Gaussian pilot density without
assuming that the unknown density is Gaussian.

The method solves the fixed-point equation
\[
t=\xi\gamma^{[l]}(t),
\qquad
\xi
=
\left(
\frac{6\sqrt{2}-3}{7}
\right)^{2/5},
\]
using
\[
z_{n+1}=\xi\gamma^{[l]}(z_n).
\]
The authors recommend $l=5$.

Algorithm 1 produces two quantities:
\[
\widehat t_p
\qquad\text{and}\qquad
\widehat t_2.
\]
The first is used to form the pilot density
\[
p(x)
=
\frac{1}{N}
\sum_{i=1}^{N}
\phi(x,X_i;\widehat t_p),
\]
while $\widehat t_2$ is an auxiliary bandwidth used later to estimate
the roughness required for the final diffusion bandwidth. Thus, the two
bandwidths have different purposes.


\subsection{Algorithm 2: diffusion estimator}

Algorithm 2 uses the pilot $p$ to construct the final estimator. For
$a=p$,
\[
\sigma^2(x)=\frac{a(x)}{p(x)}=1.
\]

The asymptotically optimal diffusion time is
\[
t^\star
=
\left[
\frac{
E_f[\sigma^{-1}(X)]
}{
2N\sqrt{\pi}\,
\|Lf\|_2^2
}
\right]^{2/5}.
\]
The difficulty is that $f$ is unknown, and hence $\|Lf\|_2^2$ is also
unknown. The paper estimates it using
\[
\widehat{\|Lf\|_2^2}
=
\|Lg(\cdot;\widehat t_2)\|_2^2.
\]
Since the diffusion equation gives
\[
Lg=\frac{\partial g}{\partial t},
\]
we numerically approximate
\[
\frac{\partial g}{\partial t}
(x;\widehat t_2)
\approx
\frac{
g(x;\widehat t_2+\varepsilon)
-
g(x;\widehat t_2)
}{
\varepsilon
}.
\]

For $a=p$, $\sigma=1$, and therefore
\[
\widehat t^\star
=
\left[
\frac{
1
}{
2N\sqrt{\pi}\,
\widehat{\|Lf\|_2^2}
}
\right]^{2/5}.
\]

Finally, the empirical distribution is evolved under the diffusion
equation from $t=0$ to $t=\widehat t^\star$. The final estimate is
\[
\boxed{
\widehat f_{\mathrm{diff}}(x)
=
g(x;\widehat t^\star).
}
\]

The complete pipeline is
\[
X_{1:N}
\longrightarrow
(\widehat t_p,\widehat t_2)
\longrightarrow
p
\longrightarrow
\widehat{\|Lf\|_2^2}
\longrightarrow
\widehat t^\star
\longrightarrow
g(\cdot;\widehat t^\star).
\]


% ============================================================
% (b) REPRODUCTION
% ============================================================

\section{Reproduction}

We reproduce the boundary-bias experiment in Figure~1 of
\citet{botev2010}. The true density is
\[
f(x)=4(1-x)^3,
\qquad 0\leq x\leq1,
\]
which is a $\mathrm{Beta}(1,4)$ density. The paper uses
\[
N=1000,
\qquad
\sqrt{t}=0.05248.
\]
Since this bandwidth is explicitly given, Algorithm~1 is not required
for this particular reproduction.

The ordinary Gaussian KDE is
\[
\widehat f_G(x;t)
=
\frac{1}{N}
\sum_{i=1}^{N}
\phi(x,X_i;t).
\]
For the boundary-aware estimator, the reflection kernel is
\[
\kappa(x,X_i;t)
=
\sum_{k=-\infty}^{\infty}
\left[
\phi(x,2k+X_i;t)
+
\phi(x,2k-X_i;t)
\right],
\]
giving
\[
\widehat f_B(x;t)
=
\frac{1}{N}
\sum_{i=1}^{N}
\kappa(x,X_i;t).
\]

The ordinary Gaussian estimator loses probability mass outside
$[0,1]$ near the boundary. The reflected estimator instead accounts
for the boundary and should therefore more closely follow the true
density near $x=0$.

\paragraph{Unspecified details and implementation choices.}
The paper specifies the target density, $N$, and bandwidth, but does not
provide the exact simulated sample or random seed. We therefore generate
a new sample from $\mathrm{Beta}(1,4)$. The infinite reflection sum must
also be truncated numerically.

% Complete this after implementation:
%
% Random seed: TODO
% Number of grid points: TODO
% Reflection truncation rule: TODO
% Numerical tolerance: TODO
% Any other choices: TODO

\begin{figure}[t]
  \centering
  % Uncomment after the result exists:
  % \includegraphics[width=0.82\linewidth]
  % {results/figure1_reproduction.pdf}
  \fbox{
    \parbox[c][3.2cm][c]{0.78\linewidth}{
      \centering
      TODO: insert Figure 1 reproduction
    }
  }
  \caption{
    Reproduction of the boundary-bias experiment.
    TODO: complete after running the experiment.
  }
  \label{fig:reproduction}
\end{figure}

\paragraph{Result.}
TODO: State whether the main qualitative result was reproduced.
Describe the behaviour of the Gaussian KDE and reflected estimator near
$x=0$, and report any important numerical differences from the paper.


% ============================================================
% (c) VARIATION
% ============================================================

\section{Variation: sample size}

Our proposed variation is the sample size $N$. We investigate how the
accuracy of the diffusion estimator changes when the amount of
available data changes. Subject to the course requirement that another
group using the same paper does not choose the same variation, we use
\[
N\in\{50,100,250,500,1000,5000\}.
\]

For each value of $N$, data are generated from the same
$\mathrm{Beta}(1,4)$ target. The complete method is then recomputed:
Algorithm~1 gives $\widehat t_p$ and $\widehat t_2$, the pilot $p$ is
constructed, Algorithm~2 gives $\widehat t^\star$, and the final
diffusion density is evaluated.

Since the true density is known, estimation error is measured by
integrated squared error,
\[
\operatorname{ISE}
=
\int
\left[
g(x;\widehat t^\star)-f(x)
\right]^2dx.
\]
For $R$ independent repetitions, we report
\[
\overline{\operatorname{ISE}}_N
=
\frac{1}{R}
\sum_{r=1}^{R}
\operatorname{ISE}_{N,r}.
\]

% Complete after implementation:
%
% Number of repetitions R: TODO
% Random-seed strategy: TODO
% PDE grid/domain: TODO
% ODE solver: TODO
% Solver tolerances: TODO
% Representation of the empirical initial condition: TODO

\begin{figure}[t]
  \centering
  % Uncomment after the result exists:
  % \includegraphics[width=0.80\linewidth]
  % {results/sample_size_experiment.pdf}
  \fbox{
    \parbox[c][3.2cm][c]{0.76\linewidth}{
      \centering
      TODO: insert sample-size experiment
    }
  }
  \caption{
    TODO: mean integrated squared error as a function of sample size.
  }
  \label{fig:variation}
\end{figure}

\paragraph{Result.}
TODO: Report the actual numerical trend after running the experiment.
State the mean error for the tested sample sizes and discuss what the
result suggests about the effect of sample size. Do not fill this section
until the experiment has been executed.


% ============================================================
% REFERENCES -- EXCLUDED FROM THE FOUR-PAGE LIMIT
% ============================================================

\bibliographystyle{plainnat}

\begin{thebibliography}{1}

\bibitem[Botev et~al.(2010)Botev, Grotowski, and Kroese]{botev2010}
Z.~I. Botev, J.~F. Grotowski, and D.~P. Kroese.
\newblock Kernel density estimation via diffusion.
\newblock \emph{The Annals of Statistics},
  38(5):2916--2957, 2010.
\newblock doi:10.1214/10-AOS799.

\end{thebibliography}


% ============================================================
% SIGNED AI STATEMENT -- EXCLUDED FROM FOUR-PAGE LIMIT
% ============================================================

\newpage
\section*{AI statement}

TODO: Paste the exact AI statement required by the course here and sign
it. Do not replace the official wording with model-generated wording.


\end{document}